
\documentclass[a4paper,12pt]{article}
\usepackage{amsmath}
\usepackage{graphicx}
\usepackage{multicol}
\usepackage{geometry}
\usepackage{fancyhdr}
\usepackage{tikz}
\usepackage{eso-pic}

\geometry{top=1.5cm, bottom=1.5cm, left=1.5cm, right=1.5cm}

\pagestyle{fancy}
\fancyhf{}
\rhead{Esc. 4-124 Reynaldo Merín}
\lhead{Electrotecnia II}
\cfoot{\thepage}

\AddToShipoutPicture*{
    \put(150,550){\rotatebox{45}{\textcolor[gray]{0.95}{EXAMEN ORIGINAL NO DIGITALIZAR}}}
}

\begin{document}

\begin{center}
    \Large \textbf{Evaluación A – Electrotecnia II} \\
    \normalsize Duración: 60 minutos \\
    Alumno/a: \underline{\hspace{6cm}} \hspace{0.5cm} Fecha: \underline{\hspace{2cm}}
\end{center}

\vspace{0.3cm}

\begin{multicols}{2}

\section*{Problemas}

\begin{enumerate}
\item Un conductor de $l = 20\,cm$ con $I = 12\,A$ está dentro de un campo de $B = 0{,}8\,T$. Calcular la fuerza ejercida sobre él.

\item Determinar la corriente necesaria para que un conductor de $l = 0{,}15\,m$ en un campo de $B = 1200\,Gs$ experimente una fuerza de $0{,}45\,N$.

\item Una bobina de $N = 400$ espiras genera un flujo de $\Phi = 1{,}2 \times 10^{-3}\,Wb$ con $I = 30\,A$. Calcular $L$.

\item Calcular la F.E.M. inducida en una bobina de $L = 0{,}006\,H$ si $\frac{di}{dt} = 0{,}2\,A/s$.

\item Una bobina tiene $r = 0{,}05\,m$, $l = 0{,}3\,m$, $N = 1500$ y núcleo de aire. Calcular $L$. ($\mu_0 = 4\pi \times 10^{-7}$)

\item Un conductor de $l = 25\,cm$ con $I = 10\,A$ está en un campo de $B = 0{,}9\,T$. Calcular la fuerza.

\item Si $L = 0{,}01\,H$ y $\frac{di}{dt} = -0{,}5\,A/s$, ¿cuál es la fem inducida?

\item Aplicando la regla de Fleming, indicar la dirección de la fuerza en un conductor con corriente hacia el este en un campo vertical descendente.

\end{enumerate}

\end{multicols}

\end{document}